\section{Preliminaries}\label{sec:prelim}

\subsection{Notation}

Throughout, $\lambda>1$ is fixed, $L=\log\lambda>0$, $v(x)=2\lambda\cos(2\pi x)$, and $H=H_{\lambda,\alpha,\theta}$ is the operator~\eqref{eq:amo}. We write $S$ for the shift $(Su)(n)=u(n+1)$, $X$ for the position operator $X\delta_j=j\delta_j$, and $\un_A$ for multiplication by the indicator of $A\subset\Z$. For a frequency $\omega\in\R$ and a phase $\varphi$ we also use $H_{\lambda,\omega,\varphi}$ with $\omega$ rational. The spectrum $\Sigma$ of $H_{\lambda,\alpha,\theta}$ does not depend on $\theta$ when $\alpha$ is irrational, since lattice translations conjugate $H_{\lambda,\alpha,\theta}$ to $H_{\lambda,\alpha,\theta+k\alpha}$, the orbit $\theta+k\alpha$ is dense, and the operator depends norm-continuously on the phase.

\emph{Transfer matrices.} For $E\in\C$ and $x\in\T$ put $T(E,x)=\begin{pmatrix}E-v(x)&-1\\1&0\end{pmatrix}$ and, for $r\ge0$,
\[
A_r(E,x)=T(E,x+(r-1)\alpha)\cdots T(E,x),\qquad A_0=I.
\]
A solution of $Hu=Eu$ satisfies $(u(r),u(r-1))^{\mathsf T}=A_r(E,\theta)(u(0),u(-1))^{\mathsf T}$. For a frequency $\omega$ other than $\alpha$ we write $A_r^{(\omega)}$. The Lyapunov exponent is $L(E)=\lim_r r^{-1}\int_\T\log\|A_r(E,x)\|\,dx=\inf_r r^{-1}\int_\T\log\|A_r(E,x)\|\,dx$. For $\lambda>1$ it equals $L=\log\lambda$ at every $E\in\Sigma$~\cite[Corollary~2]{BourgainJitomirskaya2002}.

\emph{Continued fractions.} Let $p_n/q_n$ be the convergents of $\alpha$ and put
\[
\delta_n=q_n\alpha-p_n,\qquad D_n=|\delta_n|^{-1},\qquad \beta_n=\frac{\log q_{n+1}}{q_n},\qquad b_n=\frac{\log D_n}{q_n}.
\]
The classical inequalities $(q_{n+1}+q_n)^{-1}<|\delta_n|<q_{n+1}^{-1}$ give
\begin{equation}\label{eq:bn}
q_{n+1}<D_n<2q_{n+1},\qquad 0<b_n-\beta_n<\frac{\log2}{q_n}.
\end{equation}
In particular $\limsup b_n=\limsup\beta_n=\beta(\alpha)$. Denominators grow at least like Fibonacci numbers, and $|\alpha-p_n/q_n|<q_n^{-2}$. When the scale $n$ is fixed we drop the index and write $q,p,\delta,b$, and $A=\lambda^q$. We also write $s=|\sin(2\pi q\theta)|$.

All constants below may depend on $\lambda$, on $\alpha$ and on parameters that are explicitly fixed. They never depend on the phase $\theta$, the scale $n$ or the time $T$, unless this is indicated. Thresholds of the form ``for all sufficiently large $n$'' are likewise uniform in the phase in Sections~\ref{sec:poly}--\ref{sec:dt}, among phases satisfying the stated sine condition at that scale. In Section~\ref{sec:upper} the phase is fixed and thresholds may depend on it.

\subsection{The phase set}

\begin{lemma}\label{lem:Ga}
For every irrational $\alpha$ the set $\Ga$ in~\eqref{eq:Ga} has full Lebesgue measure.
\end{lemma}

\begin{proof}
The map $\theta\mapsto q_n\theta$ preserves Lebesgue measure on $\T$, and $\{x\in\T:|\sin 2\pi x|<u\}$ has measure at most $u$. Hence $\{\theta:|\sin 2\pi q_n\theta|<q_n^{-2}\}$ has measure at most $q_n^{-2}$. Since $\sum_nq_n^{-2}<\infty$, the Borel--Cantelli lemma shows that almost every $\theta$ lies in only finitely many of these sets.
\end{proof}

\subsection{Moments}

\begin{lemma}\label{lem:moments}
\begin{enumerate}
\item For every $p>0$ there is $C_p$, depending only on $p$, such that $M_p(T)\le C_p(1+T)^p$ for all $T>0$ and every bounded potential.
\item For $0<a<1$ and $T\ge1$, $M_p(T)\le T^{ap}+\bigl(M_{2p}(T)\,P(T^a,T)\bigr)^{1/2}$.
\item If $0<a<1$ and $P(T^a,T)=O(T^{-m})$ for every $m>0$, then $\au\le a$.
\end{enumerate}
\end{lemma}

\begin{proof}
1. Let $k\ge1$ be an integer. We have $[X,H]=S^{-1}-S$ and $[X^k,H]=\sum_{i=0}^{k-1}X^{i}[X,H]X^{k-1-i}$. Moving the shifts to the right with $S^{\pm1}X=(X\pm1)S^{\pm1}$ writes $[X^k,H]$ as $\sum_{\pm}c_\pm(X)S^{\pm1}$, where $c_\pm$ are polynomials of degree at most $k-1$ with coefficients bounded by a constant depending only on $k$. Since $X\delta_0=0$, Duhamel's formula $X^k\psi_t=-i\int_0^te^{-i(t-s)H}[X^k,H]\psi_s\,ds$ and induction on $k$ give $\|X^k\psi_t\|\le C_k(1+|t|)^k$. For real $p>0$ choose an integer $k$ with $2k\ge p$. Jensen's inequality for the probability distribution $|\psi_t(j)|^2$ gives
\[
\sum_j|j|^p|\psi_t(j)|^2\le\Bigl(\sum_j|j|^{2k}|\psi_t(j)|^2\Bigr)^{p/(2k)}=\|X^k\psi_t\|^{p/k}\le C(1+|t|)^p.
\] Averaging against $(2/T)e^{-2t/T}$ proves the claim.

2. Split $M_p(T)$ into $|j|\le T^a$ and $|j|>T^a$. The first part is at most $T^{ap}$, because $\sum_ja(j,T)=1$. Apply the Cauchy--Schwarz inequality to the second part, for the probability measure $a(\cdot,T)$ on $\Z$.

3. By parts 1 and 2, $M_p(T)\le T^{ap}+C(1+T)^{p}T^{-m/2}$. Taking $m\ge2p$ gives $M_p(T)\le CT^{ap}$, so $\beta^+(p)\le a$ for every $p$.
\end{proof}

\subsection{Uniform growth near the spectrum}

\begin{lemma}\label{lem:growth}
For every $h>L$ there are a bounded open set $U_h\supset\Sigma$ and a constant $C_h$ such that
\[
\|A_r(E,x)\|\le C_he^{hr}\qquad\text{for all real }E\in U_h,\ x\in\T,\ r\ge0.
\]
\end{lemma}

\begin{proof}
Fix $E_0\in\Sigma$ and $h_1=(L+h)/2$. Since $L(E_0)=L<h_1$, there is $r_0$ with
\[
\frac1{r_0}\int_\T\log\|A_{r_0}(E_0,x)\|\,dx<h_1.
\] The function $F_E(x)=\log\|A_{r_0}(E,x)\|$ is jointly continuous, so the same strict inequality holds for $E$ in a compact neighborhood $N_0$ of $E_0$. Rotation by $r_0\alpha$ is uniquely ergodic, so Birkhoff averages of the continuous functions $F_E$ converge to $\int F_E$ uniformly in $x$; by equicontinuity of $\{F_E\}_{E\in N_0}$ they converge uniformly in $E\in N_0$ as well. Hence there is $m_0$ with $\sum_{k<m}F_E(x+kr_0\alpha)\le mr_0h$ for all $m\ge m_0$, $x\in\T$ and $E\in N_0$. Writing a product of length $r=mr_0+s$ with $0\le s<r_0$ as $m$ blocks of length $r_0$ and a remainder, submultiplicativity gives $\|A_r(E,x)\|\le C e^{hr}$ on $N_0$. A finite cover of the compact set $\Sigma$ by the interiors of such neighborhoods gives $U_h$ and $C_h$.
\end{proof}

\subsection{Rational spectra}

For a rational $p/q$ in lowest terms let $\Sigma_{p/q}=\bigcup_{\varphi\in\T}\sigma(H_{\lambda,p/q,\varphi})$.

\begin{lemma}\label{lem:window}
Let $H_0$ be a bounded Schr\"odinger operator on $\ell^2(\Z)$ with hopping one and real potential, and let $E\in\sigma(H_0)$. For every integer $M\ge1$ there is a unit vector $u$ supported in an interval of length $2M$ with $\|(H_0-E)u\|\le 8/M$.
\end{lemma}

\begin{proof}
By Weyl's criterion there is a finitely supported vector $w$ with $\|(H_0-E)w\|\le M^{-1}\|w\|$. For $k\in\Z$ let $\chi_k(n)=\sin\bigl(\pi(n-kM)/(2M)\bigr)$ for $kM\le n\le (k+2)M$ and $\chi_k(n)=0$ otherwise. Then $\sum_k\chi_k(n)^2=1$, because on $[kM,(k+1)M]$ the two windows $\chi_k$ and $\chi_{k-1}$ are a sine and a cosine of the same angle. Also $|\chi_k(n+1)-\chi_k(n)|\le\pi/(2M)$, and for each $n$ at most three indices $k$ have $\chi_k(n)$ or $\chi_k(n+1)$ nonzero. Since $([S,\chi_k]w)(n)=(\chi_k(n+1)-\chi_k(n))w(n+1)$, this gives $\sum_k\|[S^{\pm1},\chi_k]w\|^2\le 3\pi^2(4M^2)^{-1}\|w\|^2$. Hence $\sum_k\|\chi_kw\|^2=\|w\|^2$ and
\[
\sum_k\|(H_0-E)\chi_kw\|^2\le2\sum_k\|\chi_k(H_0-E)w\|^2+2\sum_k\|[H_0,\chi_k]w\|^2\le\frac{2+6\pi^2}{M^2}\,\|w\|^2.
\]
Some $k$ therefore has $\|(H_0-E)\chi_kw\|\le (8/M)\|\chi_kw\|$.
\end{proof}

\begin{lemma}\label{lem:ratspec}
There is $C$ such that for every rational $p/q$:
\begin{enumerate}
\item every point of $\Sigma_{p/q}$ is within $C|\alpha-p/q|^{1/2}$ of $\Sigma$, and every point of $\Sigma$ is within $C|\alpha-p/q|^{1/2}$ of $\Sigma_{p/q}$;
\item for every $\varphi$, every point of $\Sigma_{p/q}$ is within $2\pi\lambda/q$ of $\sigma(H_{\lambda,p/q,\varphi})$.
\end{enumerate}
\end{lemma}

\begin{proof}
1. Let $E\in\sigma(H_{\lambda,p/q,\varphi})$ and take $u$ from Lemma~\ref{lem:window}, supported in $[c-M,c+M]$. Choose $\theta'$ with $\theta'+c\alpha=\varphi+cp/q$. For $|n-c|\le M$ the potentials of $H_{\lambda,\alpha,\theta'}$ and $H_{\lambda,p/q,\varphi}$ differ by at most $4\pi\lambda M|\alpha-p/q|$. Hence $\dist(E,\Sigma)\le 8/M+4\pi\lambda M|\alpha-p/q|$, and $M=\lceil|\alpha-p/q|^{-1/2}\rceil$ gives the first claim. The second claim is proved in the same way with the roles exchanged.

2. Translation by $k$ sites conjugates $H_{\lambda,p/q,\varphi}$ to $H_{\lambda,p/q,\varphi+kp/q}$, and $\{kp/q\}$ is the full grid $q^{-1}\Z/\Z$. Every phase $\varphi'$ is within $1/(2q)$ of a grid translate of $\varphi$, and the potential is $4\pi\lambda$-Lipschitz in the phase.
\end{proof}
